% Generated draft: requires mathematical and source review.
\documentclass[11pt]{article}
\usepackage[margin=25mm]{geometry}
\usepackage{fontspec}
\setmainfont{Latin Modern Roman}
\usepackage{amsmath,amssymb,mathtools,unicode-math}
\setmathfont{Latin Modern Math}
\usepackage{enumitem,needspace,xurl,hyperref}
\setlength{\parindent}{0pt}
\setlength{\parskip}{5pt}
\setlength{\emergencystretch}{3em}
\newcommand{\N}{\mathbb{N}}
\newcommand{\Z}{\mathbb{Z}}
\newcommand{\Q}{\mathbb{Q}}
\newcommand{\R}{\mathbb{R}}
\newcommand{\C}{\mathbb{C}}
\newcommand{\F}{\mathbb{F}}
\newcommand{\A}{\mathbb{A}}
\newcommand{\PP}{\mathbb{P}}
\newcommand{\E}{\mathbb{E}}
\newcommand{\eps}{\varepsilon}
\newcommand{\dd}{\,\mathrm{d}}
\begin{document}
\section*{UnsolvedMath ID: 30000074}
\textbf{Source problem code:} catalogue entry
\subsection{Definitions and mathematical statement}
\textbf{Two requested definitions.} Give a correct definition of the category \(D(X,B)\) and an intrinsic definition of perverse point sheaves. The source does not define \(D(X,B)\); that missing definition is part of the problem, not notation to be replaced automatically by \(D^b(\operatorname{Coh}X)\). The attached literature assessment places the request in the context of Kawamata log terminal (KLT) pairs, but does not specify the intended field, dimension, boundary coefficients, or axioms for the category.\par In the usual characteristic-zero birational setting, a KLT pair consists of a normal variety \(X\) and a boundary divisor \(B\), with \(K_X+B\) real Cartier and every log discrepancy positive. For a log resolution \(f:Y\to X\), one writes \(K_Y=f^*(K_X+B)+\sum_E a(E,X,B)E\), including the relevant strict transforms in this discrepancy notation, and requires \(a(E,X,B)>-1\) for all divisors over \(X\). A rational-boundary version uses rational coefficients and \(\Q\)-Cartier \(K_X+B\). These are explanatory conventions; the record does not determine which coefficient setting the desired category must handle.\par The first request is to attach the intended category to such a pair with a precise construction and appropriate independence of any auxiliary presentation. The required relationship with ordinary derived categories, stacks, or birational maps is not specified, so no particular candidate or equivalence criterion is asserted as the answer. The second request asks for a definition of perverse point sheaves in terms of the intended intrinsic geometric or categorical data. A perverse coherent sheaf is an object in a specified perverse coherent heart of a derived category, and the point objects require further defining conditions. In relative constructions that heart may depend on a morphism \(X\to Y\); such a morphism, its hypotheses, the perversity, and what \emph{intrinsic} permits as input must be specified before this becomes a determined definition problem. A relative construction must not silently be declared intrinsic to \((X,B)\) alone.\par The supplied literature triage, checked 2026-08-21, labels the problem partially solved. It reports relative perverse coherent and perverse point sheaf constructions in important threefold-flop settings, including canonical covering stacks, while cautioning that these results do not establish a general \(D(X,B)\) theory for arbitrary KLT pairs. This is a dated assessment requiring independent verification, not a claim of a completed general definition.
\subsection{Short English statement}
Define a category attached to the pair \((X,B)\) and characterize perverse point sheaves intrinsically. The supplied triage reports relative constructions for certain threefold flops, while leaving the general KLT-pair theory unresolved. The category's required properties and the allowed intrinsic input are not specified.
\subsection{Sources}
\begin{itemize}\item Local UnsolvedMath registry, permanent ID 30000074, code catalogue entry; \emph{Perverse Point Sheaves and the Category \(D(X,B)\)}, Oberwolfach Reports Open Problems.\item \emph{Classical Algebraic Geometry (2005)}, Oberwolfach Reports, DOI 10.4171/owr/2004/31, pp. 1657--1702. \url{https://doi.org/10.4171/owr/2004/31}. Proposed year and year tag: 2004; no proposer named.\item Supplied literature label \emph{Bridgeland, Flops and derived categories}: \url{https://arxiv.org/abs/math/0009053}.\item Supplied literature label \emph{Abramovich--Chen, perverse point sheaves on stacks}: \url{https://arxiv.org/abs/math/0304354}.\item Supplied literature label \emph{Original Oberwolfach report}: \url{https://doi.org/10.4171/owr/2004/31}.\item Literature triage checked 2026-08-21, status \texttt{partially\_solved}; record updated 2026-09-16. Verify dated status and cited claims independently.\end{itemize}
\end{document}
